\chapter{The positive-$n$ thin R10 layer: an exact extra row and a bounded conductor perturbation}
\label{component:pef37475d831e}
\noindent\textbf{Component }\path{pef37475d831e}\par
{\small\noindent Source: \path{R16/proofs/r10/r10_thin_layer_contact.tex}\par}
{\small\noindent SHA256: \path{ef37475d831e144531c59d8d80c245a3f315a6c02fdf0db156d80f66e82e481c}\par}
{\small\noindent Inventory use: TL1-TL10\par}
\medskip
\section{The positive-$n$ thin R10 layer: an exact extra row and a bounded conductor perturbation}

\paragraph{Domain.}
All polynomial identities below use the convention $F/4$, as in
R10.9; in particular the original $n$ packet is $nD_uD_vX^2$.
All parent and packet identities in this section use the common normalization $F/4$.
Keep every original R10 condition and assume
$\rho=n/e^2>0$ infinitesimal. Use $s,w,b,x,B$ from R10.2.
Dividing that exact identity by $w$ gives
\[
 Q_*-(b/w)K_*=2r\Delta_*\rho w,
 \qquad b/w\sim1/t.
\tag{TL1}
\]
Thus $\psi=w/b$ is a positive bounded unit. No testing rate occurs.
The formal native C1 map OP4 is
\[
 x_1=x,\quad\eta_1=t,\quad\zeta_1=s,\quad r_1=h,\quad
 \theta_1=e,\quad t_1=\theta_1/\zeta_1=w.
\]
Write $P_d=P_n/B$. OP3--OP5 give the exact slack
\[
 s_n=2r\Delta_*\rho w/P_d,\quad
 \beta_1=-s_n/2,\quad \pi=1-s_n=1+2\beta_1>0,\quad
 \kappa=(1+2n)/\pi.
\tag{TL2}
\]
Hence $s_n/w=O(\rho)$, $\beta_1$ is infinitesimal, and
$\kappa\sim1$. The formal C1 parent is used through its explicit
identities, not imported as a physical C1 source theorem.

\paragraph{Exact coefficient comparison.}
Let $F_1$ be the native C1 parent at this same contact, normalized by
$d_1=1$, and $F_2$ the actual C2 parent, with $d_2=1$.
For a coefficient $L$ put $\delta L=L_2-\kappa L_1$.
Set $K=1+x$, $c_a=1+xr$, $S_c=1+B(1-h)+xr$, $Q_c=1-hB$,
and $\Theta=n/B=\rho w\psi$. The following short polynomials are
bounded with the displayed unit denominators:
\[
 P_\kappa=c_a^2+sx(1+t)+sx^2(r+he+t),
\]
\[
 P_D=tx^2r^2+x\{t-1+e(1-h-2t)\}-1,\quad
 P_\alpha=1+s-2e+x\{r^2+s(r+he)\}.
\]
Exact cancellation of the original C2 and native C1 formulas gives
\[
\begin{array}{c|l}
 \delta m&0\\
 \delta C&-\Theta c_a/(1-h)\\
 \delta(Ah)&-\Theta t c_aS_c/[K(h+t)]\\
 \delta D&\Theta S_cP_D/[K(h+t)tj]\\
 \delta\alpha&-nxP_\alpha/(KQ_c)\\
 \delta d=1-\kappa&-2\Theta P_\kappa/[K(h+t)]\\
 \delta\beta+\delta d/2&-n.
\end{array}
\tag{TL3}
\]
The new checker verifies all seven identities and TL2 exactly.
In particular $\delta C,\delta D,\delta d,\delta(Ah)=O(\rho w)$,
while $\delta\alpha=O(nx)$. There is no surviving $1/h$ after the
large $A$ event is centered.

\paragraph{The formal C1 local identities still supply a strict base.}
The native $\eta_1$ is a unit, $k_1,\zeta_1,r_1$ are infinitesimal,
and $e_{1,\mathrm{scale}}^2=tb$ is bounded. Separate its positive-unit
and infinitesimal types. Use respectively the exact centered-$A$
CA5/CA8 construction (or its CS1 scaling) and the SE4/SE9--SE11
construction. This use concerns the printed formulas and conductor
lemmas; the physical-source theorem is not being invoked at $\beta_1<0$.
Here are all changes to their source premises:

The formulas and native stationary identities are rational identities,
valid at the formal OP4 contact. All native gaps remain positive; $M=g$
and $J=j$ are units. The other native event coefficients are positive,
with $C_1,D_1,A_1h$ positive units and $A_1$ unbounded. The opening
ratio $t_1/e_{1,\mathrm{scale}}^2=w/(tb)$ is a positive bounded unit
by TL1, which supplies the only scale bound previously obtained from
$\beta_1\ge0$. All four retained packets SD3 involve the positive
paired coefficient $\beta_1+1/2=\pi/2$, never a standalone
$\beta_1$ square. In SE4, the coefficient of the reserved normal packet
is $(2+4\beta_1)\nu$, so its subtraction of $\nu$ is still positive,
since $\beta_1$ is infinitesimal. The remaining source bounds, finite
coquotients and strict residue matrices are unchanged. A positive-unit
normalization to $C_1=2$, when used by those formulas, changes none of
these bounds. Consequently their exact local construction yields the
same strict bounded eight-root Gram and positive retained packets for
$\kappa F_1$ on its contact chart.

\paragraph{A new positive contact-zero $n$ row.}
Use the inherited pure-testing coordinate and normals
\[
 \mathfrak w=(W+uY)/D_u,\quad A_t=v\mathfrak w,\quad
 X=N+pD_uA_t-uZ,\quad p=t_1/R_1,
\]
with $D_u=1+u^2,D_v=1+v^2$. In the exact CA chart the contact is
\[
 (u,\mathfrak w,v)=(k,\delta k,\delta k^2E),\quad
 K_c=\delta^2Ek^2=\zeta_1K/Q.
\]
Here $Q,E$ denote the inherited CA source letters, not the angular
weight of R10. The SE chart has $k=\varepsilon j$,
$\delta=1/\varepsilon$, so $K_c=j^2E$.
Put $\gamma=r/K_c$ and retain the entire state-linear square
\[
 \boxed{\quad \mathcal N^2=nD_uD_v(X+\gamma A_t)^2.\quad}
\tag{TL4}
\]
At the actual contact $A_t=K_ck$ and $X=-rk$, so its value and
all first derivatives vanish exactly. The square root of $D_uD_v$ is
pure testing; every other radical or inverse is pure source or pure
testing. The following entire-parent identity includes its full normal:
\[
\begin{split}
 nD_uD_vX^2-\mathcal N^2
 ={}&2n\gamma D_uD_vuA_tZ
 -(2n\gamma pD_u+n\gamma^2)D_uD_vA_t^2\\
 &-2n\gamma D_uD_vA_tN.
\end{split}
\tag{TL5}
\]
There is no remaining hard quadratic from this subtraction. This row
is necessary: without it the scalar $nu^2$ lies outside the old raw
square ideal, even if $n/k^4$ is infinitesimal.

\paragraph{Exact bounded stress after the positive packets.}
Let $A_{*,1}$ be the inherited bounded centered-$A$ cap. In the actual
parent retain $(A_2-\kappa A_{*,1})D_uD_v(Y-huW)^2$, which is positive
because $A_2$ is unbounded and $\kappa A_{*,1}$ bounded. Comparing with
the corresponding retained packet of $\kappa F_1$ replaces the whole
$A$ difference by
\[
 \delta(Ah)D_uD_v(2uWY-hu^2W^2).
\]
After retaining TL4, the full difference of the bounded remainders is
exactly
\[
\begin{split}
 \mathscr E={}&\bigl[nD_uD_vX^2-\mathcal N^2\bigr]
   -nD_v(X+uZ)^2+\delta C D_v(uW-Y)^2\\
 &+(\delta D+\delta d/2)R_{D,\mathrm{event}}^2
   +\delta\alpha D_u(vY-Z)^2\\
 &+(\delta d/2)D_u\{(vY-Z)^2-D_vZ^2\}\\
 &+\delta(Ah)D_uD_v(2uWY-hu^2W^2),
\end{split}
\tag{TL6}
\]
where $R_{D,\mathrm{event}}=uW+uvX-Y-vZ$ is the original event
row. This follows from TL3; in particular the paired beta difference
is exactly $-n$. Every scalar, first-normal and hard-normal term is
present. The full expression and its first derivatives vanish at the
unchanged source contact, because both full parents and every retained
square do.

\paragraph{Finite support and bounds for the unchanged conductor.}
Work on the local chart $Z=1$, with $u,v,\mathfrak w$ positive
infinitesimal. At zero normals set
\[
 V_0=1-puv\mathfrak w,\quad q_B=D_uvV_0,\quad
 uW=u\mathfrak w+u^2h_1vV_0,\quad
 Y=u\mathfrak w-h_1vV_0,\quad h_1=a_1/T_1.
\tag{TL7}
\]
These are exact inherited state inverses. The rows $uW-Y$ and
$R_{D,\mathrm{event}}$ are in $(v)$ at zero normals;
$X+u=pD_uA_t$. The centered $A$ expression in TL6 is a quadratic
in $(v,u\mathfrak w)$ with bounded coefficients. The signed-core
line has no constant: $(vY-1)^2-D_v$ is explicitly
$(vY)^2-2vY-v^2$. The sole state-constant term is
$\delta\alpha D_u$. Finally TL5 is a combination of
$uA_t=(v)(u\mathfrak w)$ and $A_t^2$, with first normal in $(A_t)$.
Thus this is an explicit raw allocation, not an ideal-membership guess.

For the unit-$e_{1,\mathrm{scale}}$ chart use
\[
 (A,B_c,C_c,D_0)
 =(\mathfrak wv,\delta v,\delta E u\mathfrak w,\delta^2k^3E).
\]
Every scalar coefficient of $\mathscr E/\rho$ in
$(A,B_c,C_c)^2+(D_0^2)$ is bounded with unit denominators. Both complete
first normals belong to $(A,B_c,C_c)$ with bounded coefficients, and
the hard-normal entries are bounded. The decisive source coquotients
are
\[
 \Theta/\rho=\delta^2QEw^2/K,\quad
 n\gamma^2/\rho=w^2r^2Q^2/K^2,\quad
 nx/(\rho D_0^2)=Q^2w^2/K^2.
\tag{TL8}
\]
The remaining $n\gamma$ coquotient has an additional $\delta^2k^2E$.
All apparent negative $k$ powers have therefore canceled before the
conductor is used.

For the infinitesimal-$e_{1,\mathrm{scale}}$ chart use the exact SE
parameters $w=t_1=\varepsilon^2\tau$,
$\zeta_1=j^2QE/K$, $k=\varepsilon j$, and
\[
 (A,B_c,C_c,D_0)
 =(\varepsilon\mathfrak wv,v,Eu\mathfrak w,\varepsilon^2j^3E).
\]
The full raw scalar allocation of $\mathscr E/\rho$ is
\[
 f_AA^2+f_{BB}B_c^2+f_{BC}B_cC_c+f_{CC}C_c^2+f_0D_0^2.
\tag{TL9}
\]
All five coefficients are bounded; after division by $\rho$ each has a
factor $\varepsilon^2$. The $A^2$ coefficient has this factor by TL10;
all its other terms carry $p=\varepsilon^2\tau/R_1$, because
$\delta m=0$. The coefficients of $B_c^2,B_cC_c,C_c^2$ have it
from $\Theta/\rho$, and the constant coefficient from the final
identity in TL10. In particular
\[
 \Theta/\rho=\varepsilon^2\tau^2QE/K,
\]
\[
 \frac{n\gamma^2}{\rho\varepsilon^2}
 =\frac{\varepsilon^2\tau^2r^2Q^2}{K^2},\quad
 \frac{n\gamma}{\rho\varepsilon}
 =\frac{\varepsilon^3j^2E\tau^2rQ^2}{K^2},\quad
 \frac{nx}{\rho D_0^2}
 =\frac{\varepsilon^2\tau^2Q^2}{K^2}.
\tag{TL10}
\]
The first-normal coefficients are in $(A,B_c,C_c)$ with bounded
arrays, also divisible by $\varepsilon^2$ after division by $\rho$;
the new $A$ coefficient has the stronger $\varepsilon^3j^2$ factor
printed in TL10. The extra term $-nD_v(X+u)^2$ has scalar coefficient
$-np^2D_u^2D_v$ on $A_t^2$, first normal
$-2npD_uD_vA_tN$, and hard term $-nD_vN^2$; TL10 and
$p=\varepsilon^2\tau/R_1$ bound each normalized coefficient.

Exactly as in SE8, all dependence of the raw coefficients on
$\mathfrak w$ comes through $V_0$ and has a factor
$p=O(\varepsilon^2)$. The additional TL5 coefficients are independent
of $\mathfrak w$. Modulo $\varepsilon$, scalar raw coefficients are
even in $u,v$; the normal raw coefficients are independent of
$\mathfrak w$. Clearing powers of $D_u$ and the source unit
denominators preserves these properties. These are precisely the
finite divisibility premises of SE7: the first scalar conductor has
$T_{\mathfrak w}\in(\varepsilon)$, the second-stage remainders have
$T_{\mathfrak w}\in(\varepsilon,v)$, and each full normal has
$T_{\mathfrak w}\in(\varepsilon)$.

For either chart, TL6 has exact contact value and first-jet zero.
Apply the inherited monic contact lift to $\mathscr E/\rho$:
CC3--CC6 in the first chart and the division-free SE7 lift in the
second. The raw support restriction retains only $D_0^2$, not arbitrary
mixed $D_0$ products. The proofs use only finite monic division and
bounded unit denominators; extending their coefficient ring by the
bounded source letters in TL8--TL10 changes nothing. They give a
bounded Gram matrix on the same eight inherited roots for
$\mathscr E/\rho$. Hence the exact stress $\mathscr E$ has an
infinitesimal Gram matrix on that list. No inverse of $\rho$ occurs in
its final coefficients: division by $\rho$ here records the common
positive source factor before a linear finite lift.

Adding this infinitesimal matrix to the strict formal-C1 base Gram
preserves its positive residue. Restore all its retained packets,
the modified positive centered-$A$ packet, and TL4. This proves the
full actual C2 parent on the thin-layer contact chart, with both
complete normals, all hard couplings and every independent rate.
The direct actual-C2 outer/state argument and the all-R10 remote
argument are required for the all-tests join; no negative formal-beta
packet is used in either join.

\paragraph{New exact receipts.}
\path{check_thin_stress.py} verifies fourteen exact source, row and scaling
identities. \path{check_thin_full_support.py} verifies five further
statements: reconstruction of the entire original-parent stress, complete
scalar restricted raw support, each full first-normal support, and absence
of higher normal degree. These are new affected checks, not reruns of the
inherited C1 suites. The finite conductor application is the analytic
argument above.
